\chapter*{Apêndice B - Questionário de Avaliação do Sistema}
\addcontentsline{toc}{chapter}{Apêndice B - Questionário de Avaliação do Sistema}
\label{apendice:questionario}

Este questionário foi aplicado de forma anônima a profissionais da área de psicopedagogia após a utilização do sistema web desenvolvido. Teve como objetivo coletar percepções sobre a facilidade de uso, a organização das informações e a adequação da plataforma às atividades de avaliação psicopedagógica.

\bigskip
\textbf{Instruções:} Para cada afirmação abaixo, indique seu grau de concordância em uma escala de 1 a 5, em que:

\begin{center}
\textbf{1} = Discordo totalmente \quad \textbf{2} = Discordo \quad \textbf{3} = Neutro \quad \textbf{4} = Concordo \quad \textbf{5} = Concordo totalmente
\end{center}

\bigskip

\begin{enumerate}

    \item A interface do sistema é de fácil compreensão e utilização.

    \bigskip
    \noindent
    \begin{tabular}{ccccc}
        \quad(\ \ ) 1 \quad & \quad(\ \ ) 2 \quad & \quad(\ \ ) 3 \quad & \quad(\ \ ) 4 \quad & \quad(\ \ ) 5 \quad
    \end{tabular}

    \bigskip

    \item As informações estão organizadas de forma clara e a navegação entre as telas é intuitiva.

    \bigskip
    \noindent
    \begin{tabular}{ccccc}
        \quad(\ \ ) 1 \quad & \quad(\ \ ) 2 \quad & \quad(\ \ ) 3 \quad & \quad(\ \ ) 4 \quad & \quad(\ \ ) 5 \quad
    \end{tabular}

    \bigskip

    \item O cadastro e a consulta dos dados dos aprendentes atendem às necessidades do processo de avaliação.

    \bigskip
    \noindent
    \begin{tabular}{ccccc}
        \quad(\ \ ) 1 \quad & \quad(\ \ ) 2 \quad & \quad(\ \ ) 3 \quad & \quad(\ \ ) 4 \quad & \quad(\ \ ) 5 \quad
    \end{tabular}

    \bigskip

    \item O gerenciamento dos instrumentos avaliativos (criação, edição e versionamento) é adequado à prática profissional.

    \bigskip
    \noindent
    \begin{tabular}{ccccc}
        \quad(\ \ ) 1 \quad & \quad(\ \ ) 2 \quad & \quad(\ \ ) 3 \quad & \quad(\ \ ) 4 \quad & \quad(\ \ ) 5 \quad
    \end{tabular}

    \bigskip

    \item O registro das aplicações e das respostas é adequado para documentar os atendimentos realizados.

    \bigskip
    \noindent
    \begin{tabular}{ccccc}
        \quad(\ \ ) 1 \quad & \quad(\ \ ) 2 \quad & \quad(\ \ ) 3 \quad & \quad(\ \ ) 4 \quad & \quad(\ \ ) 5 \quad
    \end{tabular}

    \bigskip

    \item A geração dos relatórios atende às necessidades de elaboração de devolutivas profissionais.

    \bigskip
    \noindent
    \begin{tabular}{ccccc}
        \quad(\ \ ) 1 \quad & \quad(\ \ ) 2 \quad & \quad(\ \ ) 3 \quad & \quad(\ \ ) 4 \quad & \quad(\ \ ) 5 \quad
    \end{tabular}

    \bigskip

    \item De forma geral, a plataforma contribui para a organização e o acompanhamento das atividades de avaliação psicopedagógica.

    \bigskip
    \noindent
    \begin{tabular}{ccccc}
        \quad(\ \ ) 1 \quad & \quad(\ \ ) 2 \quad & \quad(\ \ ) 3 \quad & \quad(\ \ ) 4 \quad & \quad(\ \ ) 5 \quad
    \end{tabular}

    \bigskip

    \item Você identificou dificuldades, limitações ou tem sugestões de melhoria para o sistema? Descreva livremente.

    \bigskip
    \noindent\rule{\textwidth}{0.4pt}

    \bigskip
    \noindent\rule{\textwidth}{0.4pt}

    \bigskip
    \noindent\rule{\textwidth}{0.4pt}

    \bigskip
    \noindent\rule{\textwidth}{0.4pt}

\end{enumerate}

\bigskip
\noindent\textit{Este questionário é anônimo. Nenhuma informação que permita identificar o respondente será coletada ou divulgada.}
